\documentclass[12pt,onecolumn,draftclsnofoot]{IEEEtran}
\usepackage[a4paper,margin=1in]{geometry}
\usepackage{times,graphicx,amsmath,amssymb,booktabs,cite,url}
\usepackage[ruled,linesnumbered]{algorithm2e}

\begin{document}

\title{Shared-State Coordination for a Multi-Agent Workflow Platform: Design and Repository Study}

\author{CSE3, Group 7: Krish (230607), Honey (230606), Dishita (230602), Naveen Singh\\
School of Engineering \& Technology, BML Munjal University}

\maketitle

\begin{abstract}
Multi-agent workflow platforms coordinate autonomous agents, tools, people, and persistent state, yet their reliability depends on more than the language model: task handoffs, permissions, memory, and output validation must also be made visible and controllable. This paper presents a repository-based design study of Multi-Agent-Flow, a web platform whose current codebase is branded MyceliumOps/Nexus AI. We examine its TypeScript and React application, Express services, PostgreSQL persistence through Prisma, OpenClaw gateway integration, process orchestration, shared Blackboard service, Model Context Protocol (MCP) integration, and structured run-output models. The study describes how the Blackboard represents decisions, contracts, artifacts, questions, and issues; how a bounded view is prepared for an agent; and how typed stage and final-output schemas make run state machine-readable. We also inspect the repository's Python guardrail and CrewAI/ADK bridge modules, distinguishing implemented code from claims of live production use. The available benchmark CSV contains 40 rows, but its generator samples preset distributions with a fixed random seed rather than collecting system executions. We therefore treat it as synthetic test data and do not report its values as measured system performance. Repository inspection supports concrete claims about architecture and configured limits, but it does not establish comparative accuracy, latency, safety effectiveness, or scale. The paper contributes a system map, a source-grounded workflow description, an audit of the evidence currently available, and a reproducible evaluation plan for future live trials.
\end{abstract}

\begin{IEEEkeywords}
Agentic design patterns, blackboard architecture, interoperability protocol, Model Context Protocol, multi-agent systems, security guardrails.
\end{IEEEkeywords}

\noindent\textbf{Overleaf Project Link:} Pending project creation; add the CSE3\_G7 project URL here.

\section{Introduction}
Autonomous agent systems driven by Large Language Models (LLMs) have rapidly evolved from isolated conversational chatbots into sophisticated collaborative Multi-Agent Systems (MAS) capable of executing multi-stage software engineering, analytical reasoning, and autonomous code synthesis workflows \cite{wang2024survey, xi2023rise, wooldridge2009introduction}. By partitioning a high-level computational objective into specialized functional personas---such as requirements architects, full-stack developers, security auditors, and quality assurance testers---multi-agent systems demonstrate substantial improvements in problem decomposition, reasoning depth, and task completion fidelity compared to monolithic zero-shot prompts \cite{hong2024metagpt, wu2023autogen}. Recent foundational frameworks, including AutoGen \cite{wu2023autogen}, MetaGPT \cite{hong2024metagpt}, CAMEL \cite{li2023camel}, and ChatDev \cite{qian2024communicative}, have demonstrated that simulated societies of artificial agents can co-author software artifacts through iterative natural language discourse.

These systems demonstrate several ways to structure agent work, but they do not remove the engineering choices faced by teams building a deployable platform: how to persist decisions, route work, control tool access, validate outputs, and expose run state to users. These concerns are coupled. A task can be correctly assigned yet fail if its worker cannot see a prior decision; a tool can return useful data yet introduce risk if the application trusts arbitrary URLs; and a completed run can be difficult to review if its output is only free-form text. We treat these as design questions rather than assume that any one coordination pattern inevitably causes races or context growth.

The repository offers a useful case for examining these concerns because it combines an interactive fleet-management product with orchestration and agent-integration modules. The code documents OpenClaw gateway communication as the principal external runtime path; it also includes MCP services and separate Python modules for CrewAI, ADK-style execution, and A2A-style envelopes. The latter are analyzed as repository components, not presumed to be a deployed cross-framework service. Security is likewise considered at the level supported by source: URL validation and prompt-pattern checks are present, but their existence alone does not establish resistance to adversarial attacks.

This distinction matters for both developers and users. A dashboard may show an agent as active while its remote gateway is unavailable; an output model may accept a field while the underlying statement remains unsupported; and a stored decision may be present but omitted from the context selected for a later stage. A paper about an agent platform should therefore describe visible mechanisms and the conditions under which they operate, while reserving claims about quality and safety for instrumented trials. We use this separation throughout the study: repository analysis answers what the source is designed to do, and an explicitly bounded evaluation plan identifies what remains to be measured.

This study asks two bounded questions about the implementation and the strength of its available evidence:
\begin{itemize}
    \item \textbf{RQ1:} How does the platform represent and expose shared workflow state, and what coordination guarantees can be established from its implementation?
    \item \textbf{RQ2:} Which reliability and security claims about orchestration, tools, and agent interoperability are supported by the repository, and what further evaluation is needed?
\end{itemize}

To answer these questions, we conduct a source-level case study of \textit{Multi-Agent-Flow}. The contributions are:
\begin{enumerate}
    \item We map the current web platform and agent orchestration layers from repository code and document their boundaries.
    \item We analyze the Blackboard data model and bounded prompt projection, including design properties and untested concurrency assumptions.
    \item We distinguish implemented interfaces and validation mechanisms from demonstrated production outcomes.
    \item We audit the benchmark artifact and specify a live evaluation protocol that would support valid comparative claims.
\end{enumerate}

The remainder of this paper reviews related work, defines the coordination problem, describes the implementation, explains the repository-review method, reports source-grounded findings, and closes with limitations and a plan for live evaluation.

\section{Related Work}
LLM-based multi-agent systems combine language-model reasoning with explicit roles, task boundaries, tools, and communication. Prior work studies both general agent interaction and software-development workflows. AutoGen describes conversable agents and flexible human or tool participation \cite{wu2023autogen}. CAMEL studies role-conditioned agent communication \cite{li2023camel}. MetaGPT structures collaboration around software-engineering artifacts and standard operating procedures \cite{hong2024metagpt}; ChatDev models development as a sequence of communicative phases \cite{qian2024communicative}. These works motivate separating responsibilities and making handoffs explicit. Their reported benchmark results concern their own systems and tasks; they should not be transferred to the platform studied here.

A second theme is how agents retain and share information. The blackboard model treats a shared workspace as a place where specialized knowledge sources publish partial results for others to use \cite{nii1986blackboard}. Modern agent systems also use conversation histories, summaries, retrieval, or structured memory. These mechanisms trade off visibility, freshness, and context size. A durable shared store can make decisions and open issues inspectable, while a prompt projection can constrain what enters a model call. Neither property guarantees conflict-free writes or correct retrieval; those require transaction and concurrency tests. The repository's Blackboard is therefore examined as an implementation choice whose intended properties must be distinguished from measured outcomes.

Reflexion uses textual feedback to update an agent's memory across attempts, illustrating that memory can be treated as part of an iterative control process rather than merely a longer transcript \cite{shinn2023reflexion}. A shared workflow board differs in purpose: it records task artifacts and coordination state that may be consumed by multiple agents and application services. That difference raises practical questions about authorship, resolution, stale entries, and which entries should be included in a prompt. A board needs a policy for retaining durable decisions while not allowing old artifacts to dominate current work. The repository's addressed-item ordering and capped projection expose such a policy in code, but a source review cannot determine whether its selected context is sufficient for real tasks.

Tool use creates a related boundary problem. ReAct interleaves reasoning and actions, making external observations part of an agent's next decision \cite{yao2023react}. Toolformer studies how models learn to invoke APIs \cite{schick2023toolformer}; Gorilla focuses on selecting and calling APIs \cite{patil2023gorilla}. The Model Context Protocol defines a common client-server interface for exposing tools and contextual resources \cite{mcp2026spec}. A shared protocol can reduce custom integration work, but it does not itself guarantee authorization, safe arguments, or trustworthy tool output. The platform's MCP code is considered as an integration surface, not as evidence of a measured improvement over bespoke calls.

Security studies show that natural-language instructions and retrieved content can create attack paths in LLM-integrated applications. Indirect prompt injection has been demonstrated against systems that combine models with external content and actions \cite{greshake2023not}; HackAPrompt examines prompt hacking across a global competition and illustrates why simple instruction-following assumptions are unsafe \cite{yu2023hackaprompt}. Llama Guard represents a different approach, using a learned classifier and an explicit risk taxonomy to classify prompts and responses \cite{inan2023llama}. Pattern matching can block known strings but cannot substitute for permission checks, isolation, and adversarial evaluation. Structured schemas constrain output shape, while application-level checks must still validate meaning and authorization.

Table \ref{tab:related-work} groups representative approaches by the design problem they address. The comparison is conceptual: the studies use different tasks, datasets, and metrics, so their scores are not directly comparable. This paper's focus is narrower. It documents how a single platform combines a persistent workflow application, a Blackboard projection, typed outputs, and gateway/tool integrations, then audits whether the repository contains evidence for stronger performance claims.

\begin{table}[t]
\centering
\caption{Representative Work and Its Relevance to This Study}
\label{tab:related-work}
\small
\begin{tabular}{@{}p{0.19\linewidth}p{0.28\linewidth}p{0.23\linewidth}p{0.22\linewidth}@{}}
\toprule
\textbf{Work} & \textbf{Problem and method} & \textbf{Evaluation context} & \textbf{Relevance and limitation} \\
\midrule
AutoGen \cite{wu2023autogen} & Flexible conversation among agents, humans, and tools & Multi-task agent demonstrations & Shows interaction patterns; not a platform-level baseline here \\
CAMEL \cite{li2023camel} & Role-conditioned dialogue between agents & Simulated task conversations & Motivates role separation; shared-state guarantees are outside its focus \\
MetaGPT \cite{hong2024metagpt} & Software tasks organized with roles and intermediate artifacts & Software-engineering benchmarks & Motivates structured handoffs; results cannot be compared without a shared task set \\
ChatDev \cite{qian2024communicative} & Software development through communicative phases & Software-generation tasks & Motivates staged workflows; does not evaluate this repository \\
Blackboard model \cite{nii1986blackboard} & Shared workspace for specialized knowledge sources & Historical AI problem-solving systems & Provides the coordination analogy; modern database behavior needs separate testing \\
ReAct \cite{yao2023react} & Alternating reasoning and tool action & Knowledge and decision tasks & Frames tool-mediated execution; not a multi-user fleet platform \\
\bottomrule
\end{tabular}
\end{table}

\section{Background and Problem Formulation}
A workflow run is a sequence of agent or human activities that transforms a user goal into artifacts and a final status. Let $x$ be the user request, $A$ the available agents, $T$ the available tools, and $B_t$ the shared state at step $t$. An orchestration system chooses a worker and the context supplied to it, then records the result as a state update. A useful system should preserve enough history for coordination while limiting irrelevant context and controlling access to external actions.

We model one step as
\begin{equation}
(a_t, c_t, \tau_t) = \pi(x, B_t, P_t), \label{eq:dispatch}
\end{equation}
where $\pi$ is the dispatch policy, $a_t \in A$ is the selected agent, $c_t$ is the context projection, $\tau_t \in T \cup \{\emptyset\}$ is an optional tool, and $P_t$ represents permissions and process configuration. After execution, the result $r_t$ is validated and appended or persisted:
\begin{equation}
B_{t+1} = U(B_t, V(r_t)), \label{eq:update}
\end{equation}
where $V$ is output validation and $U$ is the state-update operation. These equations describe a design abstraction, not a claim that the implementation guarantees atomicity.

The repository contains persistent process and task models, a Blackboard entry model, a service that creates a bounded agent-facing view, structured stage and final outputs, API routes, and gateway/tool adapters. The input is a user request and process configuration; the expected output is a run record containing status, stage summaries, deliverables, issues, and metrics. Assumptions include a configured database, authorized user, available gateway or model provider, and agents that can read and post in the workspace. The source review does not verify these assumptions in a live deployment.

The study is out of scope for proving model quality, production availability, security against all attacks, or superiority over another orchestration system. It does not treat code comments, planned capabilities, seeded synthetic data, or fallback strings as successful production executions. These boundaries determine the claims made in Sections V and VI.

\section{Proposed Methodology}
This section describes the design visible in the repository. Fig. \ref{fig:workflow} shows a typical path from user request to stored run output. The actual path depends on configuration: an external gateway or provider may be unavailable, and not every optional module participates in every run.

\begin{figure}[t]
\centering
\fbox{\begin{minipage}{0.94\linewidth}
\centering
\fbox{User request and process configuration}
$\rightarrow$
\fbox{Authenticated API and validation}
$\rightarrow$
\fbox{Orchestration plan and agent execution}\\[6pt]
$\downarrow$\\[2pt]
\fbox{Blackboard context and tool/gateway calls}
$\rightarrow$
\fbox{Stage output validation}
$\rightarrow$
\fbox{Persist run state and notify client}
\end{minipage}}
\caption{Source-derived workflow from request intake through persistence and client updates.}
\label{fig:workflow}
\end{figure}

Fig. \ref{fig:workflow} highlights the system boundaries that matter for reproducibility: request validation, execution dependencies, shared context, output validation, and persistence. A live trace should record each transition and failure so that a future experiment can distinguish model errors from infrastructure errors.

\subsection{Application and Orchestration Layers}
The frontend is a React and TypeScript application with pages for agents, boards, processes, approvals, metrics, and settings. It communicates with an Express backend. The backend separates routes, middleware, and services; request validation and authentication are implemented as middleware, while Prisma models and services provide persistence. Socket.IO is used for real-time event delivery. The OpenClaw client provides a WebSocket RPC boundary to remote gateways. These boundaries help isolate UI, domain logic, storage, and external runtime behavior, but the source layout alone does not establish deployment reliability.

A process can be represented as ordered or dependent assignments, with orchestration services responsible for planning and executing work. The repository includes scheduling, task, process, approval, incident, and run-related services. Human approval is relevant where an agent action has consequences beyond generating text. A robust deployment should keep approval state, actor identity, and the exact proposed action in the audit trail.

The platform also has to manage boundaries between tenants and between people and agents. The backend includes organization membership and role checks, API-key authentication, agent tokens, request validation, and resource-access middleware. These mechanisms shape which user can initiate a run and which records can be accessed. Their effectiveness depends on consistent enforcement at every route and service boundary, as well as correct database scoping. This study notes the mechanisms as architectural elements but does not claim that every authorization path is free of defects, because that would require a dedicated threat model and tests covering cross-tenant access.

\subsection{Blackboard Shared State}
The Blackboard service uses entries with kinds such as decision, contract, artifact, issue, question, answer, and note. Entries can be addressed to a role and can resolve earlier actionable entries. The service reads persisted entries, renders Markdown and JSON views, and initializes per-agent JSONL post files. The code comments describe the runner as the database writer that validates agent posts and stores them. This staged flow separates an agent's workspace write from durable database state.

The service constructs an agent-facing projection by prioritizing open items addressed to that agent, then other open items, decisions and contracts, artifacts, and remaining entries. The prompt projection has a configured maximum of 7,000 characters; a database query is limited to 500 entries. These bounds expose an explicit context policy. They do not prove that every important item fits, that no two workers overwrite the same post file, or that the database and filesystem cannot diverge. Those are testable implementation properties.

\subsection{Typed Outputs, Tools, and Guardrails}
The Pydantic models define stage and final output shapes. A stage includes status, summary, deliverables, decisions, issues, optional test metadata, tools used, and next steps. A final output combines stages, open issues, metrics, a run identifier, and messaging or orchestration metadata. Validators normalize common status synonyms, remove duplicate deliverable paths, and reject inconsistent combinations such as a completed run containing a failed stage. These checks constrain representation; they do not independently establish that a claimed test actually ran or that a deliverable is correct.

MCP routes and services provide integration points for registering servers and tools. The Python repository also contains calculator, file, web, project-data, and forecast tool modules. Tool permissioning, network policy, argument validation, and result verification must be applied at the point where the tool is called. The URL validator and prompt-pattern guardrail code are useful defense components, but static presence is not a security evaluation. In particular, keyword rules can miss paraphrases and may flag harmless text.

The separate Python framework modules include a CrewAI wrapper, an ADK-style worker wrapper, and an A2A-like JSON envelope. Their code contains fallback paths and illustrative results. This study treats them as prototypes included in the repository rather than the principal live orchestration path. The production application README instead describes OpenClaw gateway integration. A future implementation report should state which runtime was actually used for each experiment.

\subsection{Execution Procedure}
Algorithm \ref{alg:orchestration} expresses a conservative run lifecycle abstracted from the services. It emphasizes validation and recording without assuming a particular model response or successful external call.

\begin{algorithm}[t]
\caption{Auditable multi-agent workflow run}
\label{alg:orchestration}
\KwIn{Request $x$, process configuration $P$, initial state $B_0$}
\KwOut{Persisted final run record $F$}
Authenticate the requester and validate $x$ and $P$\;
Create a run identifier and persist the initial status\;
Load the permitted process context and construct an agent view $c_t$\;
\ForEach{planned stage $t$}{
  Select the configured agent and permitted tools\;
  Execute the stage through the configured runtime\;
  Validate the returned stage against the output schema\;
  Record deliverables, issues, tool calls, and status\;
  Update shared state and refresh the next context projection\;
  \If{the stage is blocked or failed}{stop or request human intervention\;}
}
Assemble and validate the final output record\;
Persist final status and emit a client update\;
\end{algorithm}

The central design choice is explicit state and schema validation. A conversation-only design is simpler for short tasks, but makes task status and unresolved issues harder for other services to query. A relational state model improves queryability and supports authorization, at the cost of migrations and consistency work across files and database records. A bounded context projection limits prompt size, but creates a relevance-selection problem: omitted context can matter. Structured outputs support predictable interfaces, but require recovery behavior for malformed model output. These tradeoffs motivate the evaluation plan rather than prove a performance advantage.

Reproducibility also depends on capturing configuration at run time. A final run record should identify the process version, agent and model configuration, tool registry, prompt or skill versions, and any human intervention. The schema contains fields for the run identifier, timestamps, orchestration pattern, status, stages, metrics, and optional preview URL; those fields provide a starting point for such a record. For a comparative study, this metadata should be immutable for each run, and the event trace should include retries and failure reasons. Without this information, a later rerun may differ because a provider model or tool changed rather than because the coordination method changed.
\section{Experimental Setup}
This paper uses a descriptive repository study rather than a live performance experiment. The study covers the checkout available on 15 November 2026. We inspected the README and setup guide, backend architecture notes, package manifests, Prisma schema and migrations, orchestration and Blackboard services, MCP routes, structured-output models, guardrail modules, and the benchmark generator and CSV. The unit and end-to-end test files were inventoried but not executed for this report. Accordingly, findings describe code structure and configured behavior; they do not claim runtime success.

The product code has a React 19 and TypeScript 5.7 frontend built with Vite, and an Express 5 and TypeScript backend. Prisma 6.19 connects the backend to PostgreSQL. Socket.IO supports real-time updates, while the OpenClaw client supplies a WebSocket RPC integration. The backend package also declares the MCP TypeScript SDK and OpenTelemetry packages. These dependencies establish available integration surfaces, not that every optional integration was exercised in a deployment.

We used source presence, explicit constants, schema constraints, and control flow as the units of analysis. A behavior was classified as implemented when a corresponding function or model exists; claims of measured quality required execution logs or a reproducible experiment. There is no captured baseline run set in the repository. The benchmark script and generated CSV were inspected to determine their provenance. Table \ref{tab:study-settings} summarizes the study boundary.

\begin{table}[t]
\centering
\caption{Repository Study Scope and Evidence Rules}
\label{tab:study-settings}
\small
\begin{tabular}{@{}p{0.28\linewidth}p{0.62\linewidth}@{}}
\toprule
\textbf{Item} & \textbf{Value or rule} \\
\midrule
Study type & Descriptive source review; no tests or live model calls run for this paper \\
Snapshot & Local checkout inspected 15 November 2026 \\
Primary platform & React 19, TypeScript 5.7, Express 5, Prisma 6.19, PostgreSQL \\
Runtime integration & OpenClaw gateway client; MCP SDK dependency; optional LLM provider services \\
Benchmarks & 40-row CSV inspected alongside its generator; treated as synthetic \\
Comparison baseline & None supported by captured, repeatable system runs \\
Evidence rule & Constants and code paths are implementation facts, not performance measurements \\
\bottomrule
\end{tabular}
\end{table}

The main metrics for a future experiment should be completion rate, schema-valid output rate, tool-call success, unresolved issues, wall-clock latency, token use, cost, and concurrent-write integrity. Each metric should be recorded from run events and measured under identical tasks, model versions, provider settings, and resource limits. A baseline should use the same application and task set with Blackboard context disabled or replaced by a documented alternative. At least three independent runs per condition are needed, with the seed and model sampling settings published. Since those data are absent, this paper does not fill in a results table with invented values.

The comparison should use a paired design: each condition receives the same prompt, process graph, tools, model endpoint, timeout, and resource envelope. Run order should be randomized or interleaved to reduce the effect of provider load and time-of-day changes. Output quality should be scored by reviewers who do not know the condition, using a rubric that defines correctness, completeness, requirement coverage, and safety separately. Report individual run records alongside mean and standard deviation, and explain missing or failed runs rather than silently excluding them. An ablation can disable one mechanism at a time, such as context projection or schema validation, while preserving the rest of the runtime. This setup would test the research questions directly; the current repository does not contain the data needed to execute that analysis.

\section{Results and Analysis}
The repository study identifies an application platform with persistent organizational data and agent-run coordination features. Fig. \ref{fig:repository-architecture} summarizes the observed boundaries. The web client communicates with the Express API; backend services combine relational state, gateway and provider integrations, and run orchestration. The diagram represents code organization, not a measured deployment topology.

\begin{figure}[t]
\centering
\fbox{\begin{minipage}{0.94\linewidth}
\centering
\textbf{React 19 / TypeScript client}\\[2pt]
Dashboard, boards, agents, processes, metrics\\[5pt]
$\updownarrow$ HTTP API and Socket.IO events\\[5pt]
\textbf{Express 5 API}\\
Authentication and organization middleware $\rightarrow$ routes $\rightarrow$ services\\[5pt]
\begin{tabular}{ccc}
\textbf{Process orchestration / Blackboard} & \textbf{OpenClaw gateway / LLM providers} & \textbf{MCP tools / observability}\\
\multicolumn{3}{c}{$\updownarrow$ Prisma data access}\\
\multicolumn{3}{c}{\textbf{PostgreSQL}}
\end{tabular}
\end{minipage}}
\caption{Source-derived view of the platform's major application and integration layers.}
\label{fig:repository-architecture}
\end{figure}

Fig. \ref{fig:repository-architecture} shows that the repository separates user-facing pages from API routes and service logic. It also shows external agent and tool integrations as distinct boundaries; the presence of an adapter does not mean the corresponding remote service was active during this study.

Table \ref{tab:implementation-evidence} summarizes the most relevant implementation evidence. The Blackboard service reads relational entries, prepares a Markdown projection, and orders items so that addressed open items and decisions can be surfaced within a prompt budget. The code sets that budget to 7,000 characters and limits a database read to 500 entries. These are configured bounds, not measurements of actual context reduction or a guarantee that concurrent writes cannot conflict. The service documentation says the runner writes database rows after agents post JSON lines; the implementation should therefore be described as a staged file-to-database path rather than as a database lock eliminating all file races.

\begin{table}[t]
\centering
\caption{Implementation Evidence and What It Supports}
\label{tab:implementation-evidence}
\small
\begin{tabular}{@{}p{0.25\linewidth}p{0.38\linewidth}p{0.28\linewidth}@{}}
\toprule
\textbf{Component} & \textbf{Repository evidence} & \textbf{Supported conclusion} \\
\midrule
Blackboard & PostgreSQL entries; JSONL agent posts; 7,000-character prompt projection & Shared-state and bounded-context mechanisms exist \\
Structured outputs & StageOutput and FinalOutput Pydantic models; schema version 1.0 & Run outputs have typed fields and consistency checks \\
Tool boundary & MCP SDK and MCP routes/services in backend & MCP integration surfaces are part of the codebase \\
Gateway boundary & OpenClaw WebSocket RPC client and gateway services & The application has an adapter for remote agent gateways \\
Guardrails & Python pattern checks and URL validation code & Some input checks exist; effectiveness is unmeasured \\
Benchmark artifact & Seeded generator creates 40 rows from preset distributions & CSV is synthetic and cannot establish system performance \\
\bottomrule
\end{tabular}
\end{table}

The structured-output models define run and stage status, deliverables, issues, test metadata, and operational metrics. A model-level consistency rule rejects a completed final run when one of its stages is marked failed; another rule prevents resolved items from appearing in the open-issues collection. This establishes validation logic in the model layer, but the rate at which an LLM response conforms to those models requires live observations.

The schema also separates a reported test from the outcome of a workflow: a stage may contain a test-run field with a boolean indicating whether a command ran and an optional pass value. This distinction is valuable because an agent-generated summary should not be confused with an independently executed test. The run model includes counts for tokens, cost, agent steps, tool calls, and Blackboard entries, but a field's existence is not evidence that production runs populated it correctly. A future measurement should validate instrumentation against provider records and server logs before computing aggregate cost or success rates.

The Blackboard's fixed prompt projection is a design parameter with two competing effects. A smaller view can reduce repeated context and make prompts easier to audit, while a strict cap can omit older decisions or long artifacts. The ordering logic prioritizes entries addressed to an agent and then other open items, but it does not provide semantic ranking. An evaluation should measure both inclusion of task-critical facts and the amount of context passed to each stage. It should also test boards with more than 500 entries and posts near the 4,000-character per-entry limit to reveal truncation and omission behavior.

The repository also includes Python modules for a CrewAI engine, an ADK-style worker, an A2A envelope bridge, and prompt-pattern guardrails. Inspection shows that these modules contain deterministic fallback behavior and illustrative outputs. Their existence is useful as prototype material, but the current code evidence does not support claims that a production CrewAI-to-ADK service passed 40 runs or that a standards-compliant A2A exchange was evaluated. Similarly, the benchmark generator sets a fixed random seed, samples values around hard-coded means and standard deviations, derives token counts from synthetic cost, and assigns confidence values from a random interval. Thus the 40 CSV records are generated data; they are not observations from the platform. The prior chart and heatmap derived from these values have been excluded from the findings.

For RQ1, the source supports the conclusion that the platform provides a shared Blackboard representation and a bounded prompt projection. It does not establish that the design outperforms peer-to-peer messaging, prevents every race, or reduces token use by a measured amount. For RQ2, the source supports the existence of MCP, gateway, schema, and guardrail components. It does not establish comparative reliability, successful cross-framework interoperability under live conditions, or a quantified reduction in security risk. The gap between implemented mechanisms and measured outcomes is the primary result of this repository study.

The file-backed posting path deserves separate evaluation from the database projection. If two agents use distinct post files, their writes may be isolated at the filename level; if names normalize to the same slug or a process retries a write, duplication or collision behavior may differ. The service's documented runner validation step is a natural place to reject malformed records and enforce per-agent limits. Tests should verify ordering, resolution, access control, and recovery when a post is only partially written. We do not infer any such pass rate from test files or comments because these tests were not executed here.

\section{Discussion}
The codebase reflects a practical systems concern: agents require durable task context, typed outputs, tool boundaries, and a user interface for reviewing work. The Blackboard makes coordination state explicit through typed entries rather than requiring every decision to be repeated as unstructured dialogue. The bounded projection provides a clear policy for selecting what an agent sees. Typed outputs give downstream services a predictable structure for status, deliverables, and issues. Together, these choices are plausible foundations for inspectable workflows, but their benefits remain hypotheses until tested against alternatives.

A meaningful next evaluation should compare Blackboard and non-Blackboard coordination on the same task set, then report both quality and cost. Concurrency testing should include workers posting to the same and different topics, interruption during posting, malformed JSON lines, duplicate resolutions, and database failure. Tool and gateway trials should log request identifiers, latency, status, and error categories. Security evaluation should use a documented attack corpus and report both missed attacks and benign inputs incorrectly blocked. These tests should be run against the deployed backend and named versions of remote models and tools.

Several limitations constrain this study. It is based on one local repository snapshot and did not run the backend, connect to a database, invoke an LLM, or inspect production telemetry. The repository includes modules from different stages of development, and descriptive documentation may be stale. The Python guardrail patterns are heuristic; strings can evade simple matching and benign text can match broad terms. URL filtering requires adversarial testing across redirects, DNS changes, IPv4/IPv6, and proxy behavior before any security claim is warranted. The synthetic benchmark is not suitable for estimating latency, cost, confidence, or comparative performance.

For practitioners, the immediate value of the implementation is operational legibility: tasks, approvals, agent status, and outputs have places in the product model and user interface. This can make it easier to supervise a fleet than to inspect an unstructured conversation alone. The corresponding responsibility is to preserve a human path for cancellation, correction, and approval, especially when agents can invoke tools or change project files. A UI that exposes status should also make stale or partial status clear; future usability work should examine whether users can identify a blocked run, inspect its evidence, and understand what an approval will authorize.

Threats to validity include selection bias in the files inspected, source behavior that may differ from deployment configuration, and provider behavior that changes over time. Future work should publish a task corpus, exact environment and dependency versions, raw event logs, run seeds, and scripts that compute aggregate values. Any claim about safety should be scoped to the tested attack set. Ethical deployment also requires human review for consequential actions, least-privilege tool credentials, tenant isolation, and a clear record of what agents changed.

\section{Conclusion and Future Work}
This paper examined the Multi-Agent-Flow repository as a source-level study of a multi-agent workflow platform. In answer to RQ1, the code represents shared coordination using Blackboard entries and constructs a prompt view capped at 7,000 characters; no comparative conflict or efficiency result can be inferred from that implementation fact. In answer to RQ2, the repository contains gateway, MCP, structured-output, and guardrail components, but it does not provide captured live evaluation data sufficient to claim measured interoperability, reliability, or security improvement.

The benchmark CSV must be treated as synthetic because its generator samples preset distributions rather than recording system executions. A repeatable evaluation is therefore the next necessary step. Future work should (1) instrument actual runs and compare coordination strategies on a fixed task set, (2) test Blackboard concurrency and recovery behavior under controlled faults, and (3) evaluate tool and prompt defenses against a published attack suite while measuring false positives. This evidence would let the authors replace architectural expectations with bounded, reproducible findings. Until then, the paper's contribution is an implementation map and an evidence audit, not a claim of measured superiority.
\section*{Code Repository and Reproducibility}
The repository reviewed for this paper is available at:\\
\url{https://github.com/krishrathi1/Multi-Agent-Flow.git}

\section*{AI-Use Declaration}
Generative AI was used to help organize and draft portions of this manuscript and summarize repository files. The authors must verify every technical statement and reference against the source material, revise the text in their own words, and disclose any additional AI use required by their course policy. No AI-generated experimental results are reported in this paper. The authors should attach the required similarity and AI-content reports; compliance with the course thresholds cannot be inferred from this declaration.

\section*{Author Contributions}
The following allocation is a provisional draft for the team to confirm before submission.
\begin{itemize}
    \item \textbf{Krish (230607):} Project coordination and backend orchestration and persistence implementation.
    \item \textbf{Honey (230606):} Agent and tool integration, including MCP-related and security components.
    \item \textbf{Dishita (230602):} Frontend and dashboard implementation, documentation, and evaluation planning.
\end{itemize}

\bibliographystyle{IEEEtran}
\bibliography{references}

\end{document}
